\documentclass[]{revtex4-2}
\usepackage{import, preamble}
\usepackage{float}

\newcommand{\aidan}[1]{\textcolor{red}{AMW: #1}}
\newcommand{\gourab}[1]{\textcolor{blue}{GG: #1}}

\begin{document}

\title{Threshold-Topology Interaction in the Global Cascade Model}
\author{Aidan Meador-Woodruff}
\email{ameadorw@ur.rochester.edu}
\affiliation{Department of Physics and Astronomy, University of Rochester, 206 Bausch and Lomb Hall, Rochester, NY 14611}
\author{Gourab Ghoshal}
\affiliation{Department of Physics and Astronomy, University of Rochester, 206 Bausch and Lomb Hall, Rochester, NY 14611}
\affiliation{Department of Computer Science, University of Rochester, Rochester, NY 14611}
\begin{abstract}
    Global cascades offer a simple graph-based model for large-scale failures in networked systems, from power grids and financial networks to neural avalanches and social contagion. While the effects of network topology and threshold distribution have been studied separately, how these two sources of heterogeneity interact remains poorly understood. We analyze a continuous two-parameter family, the beta distribution for thresholds and the beta-binomial for degrees, that interpolates between power-law and Poisson-like topologies and between heavy-tailed and nearly homogeneous thresholds. The cascade phase diagram is organized by a boundary at $\alpha_{\rm th}=1$, independent of the degree distribution and of the detailed shape of the threshold distribution, that separates a regime where cascades persist at arbitrarily high connectivity from regimes where they are eventually suppressed. Within the suppressible regime, a second boundary $\alpha_{\rm th}=\alpha_d+2$ separates threshold-dominated from degree-dominated cascades. The interplay of threshold and degree skewness, not the heaviness of the tails, is the fundamental mechanism governing cascade dynamics.
\end{abstract}

\maketitle

\section{Introduction}

Cascading failures in power grids, financial systems, and social networks are large, rare events triggered by small initial shocks. In power grids, a handful of failing components can trigger blackouts affecting millions~\citep{MotterLai2002, Yang2017} and cascade through interdependent infrastructure~\citep{Brummitt2012}. The default of a single financial institution can propagate through interbank liabilities and destabilize an entire economy~\citep{GaiKapadia2010}, a small seed of adopters can flip a social norm~\citep{Centola2010}, and a brief stimulus to a neural circuit can ignite an avalanche that recruits a large fraction of neurons~\citep{BeggsPlenz2003, ShewPlenz2013}. In epidemiology, highly right-skewed distributions of individual infectiousness produce superspreading events in which a small fraction of cases causes most transmission~\citep{LloydSmith2005}. Across these settings, agents on a network activate when the fraction of their active neighbors crosses a threshold, and newly active agents can push their neighbors past threshold. The dynamics are governed by two ingredients, the degree distribution and the threshold distribution.

In the standard model~\citep{Watts2002}, which extends Granovetter's threshold model of collective behavior to sparse random graphs~\citep{Granovetter1978}, a random network is drawn with degree distribution $p_k$, and each node is assigned a threshold $\phi$ drawn from a distribution $F$; a node of degree $k$ activates when at least $\lceil\phi k\rceil$ of its neighbors are active. A global cascade recruits a finite fraction of the network and requires a branching ratio $\kappa>1$, where $\kappa$ is the expected number of further activations triggered by a newly active node. This condition holds for arbitrary degree and threshold distributions, and Watts already treated non-constant thresholds. A constant threshold on ER graphs yields the familiar cascade window in mean degree $z$, bounded below by connectivity and above by density, while broader threshold distributions widen this window~\citep{Watts2002}. Gleeson and Cahalane extended the analysis beyond the onset condition, deriving a recursion for the final cascade fraction $\rho$ that holds for arbitrary degree and threshold distributions and finite seed fractions $\rho_0$~\citep{GleesonCahalane2007}.

Topology and thresholds have since been studied largely in isolation. On ER graphs the branching ratio $\langle k(k-1)\rangle/\langle k\rangle = z$ grows without bound, so topology alone never prevents cascades. Power-law degree distributions $p_k \propto k^{-\gamma_d}$ behave differently. Their branching ratio decreases with $\gamma_d$ and crosses unity at $\gamma_d \approx 3.48$, so such networks cannot cascade under any threshold~\citep{AielloChungLu2000}. Right-skewed topologies of this kind are also robust to random node failure, though fragile under targeted attack~\citep{Albert2000}. Threshold heterogeneity acts in the opposite direction. Broadening the threshold distribution, whether Gaussian or power-law, places more probability near zero and widens the cascade window~\citep{Watts2002, GleesonCahalane2007}.

It remains unclear how the two sources of heterogeneity interact when both are present and continuously tunable. We address this with a doubly tunable model in which both the threshold and the degree distribution interpolate continuously between homogeneous and heterogeneous extremes, a beta distribution of thresholds whose shape parameter $\alpha_{\rm th}$ controls the near-zero density and a beta-binomial degree distribution whose shape parameter $\alpha_d$ controls the right skew. The beta distribution is a well-motivated model of threshold heterogeneity. Recent work solved Granovetter's threshold model exactly on a fully connected population and showed that the near-zero shape exponent $\alpha$ alone separates the inactive from the mobilized phase at a transition $\alpha=1$~\citep{Fontanari2026}, the same exponent that sets the $\alpha_{\rm th}=1$ boundary here. The four classical limiting cases, power-law or ER degree crossed with power-law or Gaussian threshold, are the corners of this space. We find that the interaction organizes cascades into three regimes separated by a boundary at $\alpha_{\rm th}=1$. Cascades are unavoidable when $\alpha_{\rm th}<1$, conditional within a finite window of $z$ when $1<\alpha_{\rm th}<\alpha_d+2$, and suppressed by the degree distribution when $\alpha_{\rm th}\ge\alpha_d+2$. All derivations are collected in Appendix~\ref{app:deriv}.


\section{The Four Limiting Cases}
\label{sec:fourcases}

Global cascades are governed by two interacting sources of heterogeneity, the distribution of degrees and the distribution of thresholds. To disentangle these effects, we consider the four limiting cases that span the extremes, power-law or Erd\H{o}s--R\'{e}nyi (ER) topology crossed with power-law or approximately Gaussian thresholds. The four corners of this space show distinct cascade behavior and motivate the continuous parameterization developed in \S\ref{sec:bbtopo} and \S\ref{sec:phase}.

The cascade condition is stated compactly in terms of the degree generating function
\begin{equation}
    G_0(x) = \sum_{k=0}^{\infty} p_k x^k,
    \label{eq:G0}
\end{equation}
whose derivatives at unity give the mean degree $G_0'(1)=\langle k\rangle$ and the branching ratio $G_0''(1)/G_0'(1)=\langle k(k-1)\rangle/\langle k\rangle$. A global cascade requires the branching ratio to exceed unity, so a newly active node triggers on average more than one further activation. The final cascade fraction $\rho$ follows from the recursion $q_{n+1}=\rho_0+(1-\rho_0)\,G(q_n)$, where $\rho_0$ is the initial seed fraction, $q$ is the probability that a randomly chosen neighbor of a random node is active, and
\begin{equation}
    G(q) = \sum_{k=1}^{\infty} \frac{k}{\langle k\rangle}\,
    \sum_{m=0}^{k-1} \binom{k-1}{m}\, q^m(1-q)^{k-1-m}\, F\!\left(\frac{m}{k}\right),
    \label{eq:G_intro}
\end{equation}
with $F$ the threshold CDF. Expanding $G$ around the trivial fixed point $q=0$ yields $G(q)\approx\kappa q$ with the cascade parameter
\begin{equation}
    \kappa = (1-\rho_0)\sum_{k=1}^{\infty} \frac{k(k-1)}{\langle k\rangle}\, p_k\, F\!\left(\frac{1}{k}\right),
    \label{eq:kappa_intro}
\end{equation}
under the self-flipping convention in which a node of degree $k$ compares its threshold to $1/k$~\citep{GleesonCahalane2007}. Equations~\eqref{eq:G_intro} and~\eqref{eq:kappa_intro} hold for arbitrary degree distributions and threshold CDFs. The four limiting cases below substitute the appropriate $p_k$ and $F$ into Eq.~\eqref{eq:kappa_intro}.

\subsection{PL degree and PL threshold}
\label{sec:pl-pl}

When both distributions are strongly skewed, with a power-law degree and a Pareto threshold, the system is maximally heterogeneous. The degree distribution is $p_k \propto k^{-\gamma_d}$ for $k\ge1$, and thresholds are drawn from a Pareto distribution $F(x) = 1 - (x/x_{\min})^{-(\gamma_{\rm th}-1)}$ rescaled to mean $\bar\phi$ as in \S\ref{sec:betath}.

The cascade parameter is
\begin{equation}
    \kappa = (1-\rho_0)\sum_{k=2}^\infty \frac{k(k-1)p_k}{\langle k\rangle}
    \left[1 - \left(\frac{\bar\phi(\gamma_{\rm th}-2)}{\gamma_{\rm th}-1}\,k\right)^{\gamma_{\rm th}-1}\right] ,
\end{equation}
with the term in brackets cut off when the argument exceeds~1, since nodes with effective threshold above~1 never flip.

Two features dominate. First, as $\gamma_{\rm th}\to2^+$, the threshold CDF approaches a step at zero. Almost every node is vulnerable, and
\begin{equation}
    \kappa\to (1-\rho_0)\left[\frac{\zeta(\gamma_d-2)}{\zeta(\gamma_d-1)}-1\right],
\end{equation}
which is finite only when $\gamma_d>3$. In this limit cascades occur for any $\bar\phi$, confirming that extreme threshold skewness overwhelms topological constraints. Second, as $\gamma_d\to2^+$, the degree distribution becomes so broad that $\langle k\rangle$ grows without bound for fixed $k_{\max}$; the branching ratio $\langle k(k-1)\rangle/\langle k\rangle$ itself diverges, and the system is always cascadable.

\subsection{PL degree and Gaussian threshold}
\label{sec:pl-gauss}

Replacing the power-law threshold with a narrow Gaussian of width $\sigma\ll\bar\phi$ isolates the effect of topological heterogeneity. The cascade condition depends only on the fraction of nodes with threshold below $1/k$. The cascade parameter takes the form
\begin{equation}
    \kappa = (1-\rho_0)\sum_{k=2}^\infty \frac{k(k-1)p_k}{\langle k\rangle}\,
    \Phi\!\left(\frac{1/k-\bar\phi}{\sigma}\right),
\end{equation}
where $\Phi$ is the standard normal CDF.

Because $\Phi((1/k-\bar\phi)/\sigma)$ decays super-exponentially in $k$ for any $\bar\phi>0$, the sum is dominated by the first few $k$. The topology therefore sets an upper bound through the branching ratio
\begin{equation}
    \kappa \le \frac{\langle k(k-1)\rangle}{\langle k\rangle}
    = \frac{\zeta(\gamma_d-2)}{\zeta(\gamma_d-1)} - 1 .
\end{equation}
For PL degree distributions, this ratio decreases monotonically with $\gamma_d$ and crosses $\kappa=1$ at $\gamma_d\approx3.48$, as shown in Fig.~\ref{fig:four-cases-kappa}. For $\gamma_d\gtrsim3.48$, no threshold distribution, however skewed, can produce a global cascade; the degree distribution alone gives a branching ratio below one. This structural bound is absent for ER graphs, where $\langle k(k-1)\rangle/\langle k\rangle = \langle k\rangle$ and can be made arbitrarily large.

\subsection{ER degree and PL threshold}
\label{sec:er-pl}

On ER graphs with $p_k = e^{-z}z^k/k!$, the degree distribution is narrowly peaked around $z$. With a power-law threshold, the cascade parameter takes the same functional form as the PL--PL case but with the Poisson generating function
\begin{equation}
    \kappa = (1-\rho_0)\,e^{-z}\sum_{k=2}^\infty \frac{z^{k-1}}{(k-2)!}
    \left[1 - \left(\frac{\bar\phi(\gamma_{\rm th}-2)}{\gamma_{\rm th}-1}\,k\right)^{\gamma_{\rm th}-1}\right] .
\end{equation}
Because $\langle k(k-1)\rangle/\langle k\rangle = z$ grows linearly with $z$, there is no topological upper bound. At high connectivity, even a narrowly peaked degree distribution can support cascades if enough nodes have low thresholds.

The $\gamma_{\rm th}\to2^+$ limit again produces cascades at any $\bar\phi$, while large $\gamma_{\rm th}$ yields nearly homogeneous thresholds and recovers Watts' upper and lower phase transition in $z$~\citep{Watts2002}.

\subsection{ER degree and Gaussian threshold}
\label{sec:er-gauss}

This is the most restrictive case, combining a homogeneous topology with a homogeneous threshold distribution. The cascade parameter is
\begin{equation}
    \kappa = (1-\rho_0)\,e^{-z}\sum_{k=2}^\infty \frac{z^{k-1}}{(k-2)!}\,
    \Phi\!\left(\frac{1/k-\bar\phi}{\sigma}\right).
\end{equation}
For a narrow Gaussian with $\sigma\to0$, $\Phi$ approaches a step function and the system reduces exactly to Watts' original model. Cascades are confined to an intermediate range of $z$; above a critical mean degree cascades stop because $1/k < \bar\phi$ for the typical node. The branching ratio $\kappa\le z$ is not restrictive, so the threshold alone closes the cascade window.

\subsection{Synthesis}

Three patterns emerge from this comparison. First, threshold heterogeneity is the stronger driver. As $\gamma_{\rm th}\to2^+$, or in the beta framework of \S\ref{sec:betath} as $\alpha_{\rm th}\to1^+$, cascades become unavoidable regardless of topology or $\bar\phi$. Second, topology imposes a structural bound on propagation through the branching ratio. PL degree distributions with $\gamma_d\gtrsim3.48$ cannot cascade under any threshold, while ER graphs have no such limit, as shown in Fig.~\ref{fig:four-cases-kappa}. Third, all four limiting cases are continuous limits of the beta--beta binomial model introduced in \S\ref{sec:bbtopo}, where $\alpha_d$ and $\alpha_{\rm th}$ interpolate continuously between the four cases, as shown in Fig.~\ref{fig:four-cases-rho}. We develop this unification in the rest of the paper.

\begin{figure}[t]
\centering
\fbox{\parbox{0.6\columnwidth}{\centering\vspace{2cm} Placeholder for Fig.~2\vspace{2cm}}}
\caption{Branching ratio $\kappa$ contours at $\kappa=1$ in the $(\langle k\rangle,\gamma_{\rm th})$ plane for the four limiting cases, power-law degree $\times$ power-law threshold in the upper left, power-law degree $\times$ Gaussian threshold in the upper right, ER degree $\times$ power-law threshold in the lower left, and ER degree $\times$ Gaussian threshold in the lower right. Dashed lines mark the topological bound $\gamma_d\approx3.48$.}
\label{fig:four-cases-kappa}
\end{figure}

\begin{figure}[t]
\centering
\fbox{\parbox{0.6\columnwidth}{\centering\vspace{2cm} Placeholder for Fig.~3\vspace{2cm}}}
\caption{Final cascade fraction $\rho$ as a function of mean degree $z$ for the four limiting cases, showing the transition from continuous to discontinuous bifurcations across the four cases.}
\label{fig:four-cases-rho}
\end{figure}
\section{The Beta Distribution of Thresholds}
\label{sec:betath}

The beta distribution provides a natural two-parameter family of threshold distributions on the unit interval. Its PDF is
\begin{equation}
    f(x; \alpha_{\rm th}, \beta_{\rm th}) = \frac{x^{\alpha_{\rm th}-1}(1-x)^{\beta_{\rm th}-1}}{B(\alpha_{\rm th}, \beta_{\rm th})},
\end{equation}
where $B(\alpha,\beta)=\Gamma(\alpha)\Gamma(\beta)/\Gamma(\alpha+\beta)$. The parameter $\alpha_{\rm th}$ controls the near-zero behavior through $f \sim x^{\alpha_{\rm th}-1}$ as $x\to0$, so $\alpha_{\rm th}<1$ produces a divergent tail at zero and many vulnerable nodes, $\alpha_{\rm th}=1$ gives a finite limit, and $\alpha_{\rm th}>1$ suppresses probability near zero. When $\alpha_{\rm th}=\beta_{\rm th}$ the distribution is symmetric about $1/2$ and converges to a Gaussian as $\alpha_{\rm th}\to\infty$; when $\alpha_{\rm th}\to0^+$ it reduces to a power law $f\sim x^{-1}$, as shown in Fig.~\ref{fig:beta-dist}.

\subsection{Decoupling mean from shape}

A fixed mean $\bar\phi$ does not determine the threshold distribution because $m \equiv \langle x\rangle = \alpha_{\rm th}/(\alpha_{\rm th}+\beta_{\rm th})$ correlates with shape. To obtain a heterogeneous distribution with a given mean, we draw from a rescaled beta. If $X \sim \text{Beta}(\alpha_{\rm th},\beta_{\rm th})$ with mean $m$, then $\Phi = (\bar\phi/m)\,X$ has mean $\bar\phi$ while preserving the shape parameters $(\alpha_{\rm th},\beta_{\rm th})$. Nodes with $\Phi>1$ are never vulnerable, consistent with the model's definition. The CDF entering the cascade equations is then
\begin{equation}
    F(x) = \min\!\left(1,\; I_{m x/\bar\phi}(\alpha_{\rm th},\beta_{\rm th})\right),
\end{equation}
where $I_u(\alpha,\beta)$ is the regularized incomplete beta function, and the truncation accounts for rescaled thresholds above unity, which are never vulnerable,
\begin{equation}
    I_u(\alpha,\beta) \equiv \frac{1}{B(\alpha,\beta)}\int_0^u t^{\alpha-1}(1-t)^{\beta-1}\,dt .
\end{equation}

\begin{figure}[t]
\centering
\fbox{\parbox{0.6\columnwidth}{\centering\vspace{2cm} Placeholder for Fig.~1\vspace{2cm}}}
\caption{Beta threshold distributions on the left for homogeneous $\mathrm{Beta}(20,20)$ and heterogeneous $\mathrm{Beta}(2,20)$, $\mathrm{Beta}(5,20)$, $\mathrm{Beta}(10,20)$, alongside beta-binomial degree distributions on the right at $\langle k\rangle=5$ for $\alpha_d=0.25,1,5,1000$ showing the progression from right-skewed to Poisson-like.}
\label{fig:beta-dist}
\end{figure}

\subsection{Cascade condition on ER graphs}

On an Erd\H{o}s-R\'{e}nyi graph with mean degree $z$, substituting the Poisson degree distribution into Eq.~\eqref{eq:kappa_intro} gives the cascade parameter
\begin{equation}
    \kappa = (1-\rho_0)\sum_{k=2}^\infty \frac{k(k-1)p_k}{z}\,
    I_{m/(k\bar\phi)}(\alpha_{\rm th},\beta_{\rm th}),
\end{equation}
where $p_k = e^{-z}z^k/k!$. The $\kappa=1$ contour in $(z,\alpha_{\rm th})$ space separates the cascading and non-cascading regimes.

In the large-degree limit $z\to\infty$, the Poisson distribution becomes sharply peaked at $k\approx z$ and the argument of the incomplete beta function vanishes. Using the near-zero expansion of the incomplete beta function, Eq.~\eqref{eq:Iu-asym} of Appendix~\ref{app:deriv}, the sum is evaluated asymptotically in Appendix~\ref{app:deriv} to give
\begin{equation}
    \kappa \sim (1-\rho_0)\,\frac{m^{\alpha_{\rm th}}}{\alpha_{\rm th} B(\alpha_{\rm th},\beta_{\rm th})}\,
    \bar\phi^{-\alpha_{\rm th}}\,z^{1-\alpha_{\rm th}} .
    \label{eq:er-asym}
\end{equation}
The sign of the exponent $1-\alpha_{\rm th}$ determines the large-$z$ behavior, as shown in Fig.~\ref{fig:beta-er-kappa}. For $\alpha_{\rm th}<1$, $\kappa$ grows with $z$ and the system is always cascadable at sufficiently high connectivity. For $\alpha_{\rm th}>1$, $\kappa$ decays and cascades are suppressed at large $z$. The boundary $\alpha_{\rm th}=1$ is critical. Here the near-zero tail of the threshold CDF is exact, $F(1/k)\approx \beta_{\rm th} m/(\bar\phi\,k)$, and $\kappa$ approaches a constant
\begin{equation}
    \kappa_\infty = (1-\rho_0)\,\frac{\beta_{\rm th}\,m}{\bar\phi},
    \label{eq:kappa-infinity}
\end{equation}
where $B(1,\beta_{\rm th})=1/\beta_{\rm th}$ and $m=1/(1+\beta_{\rm th})$ at $\alpha_{\rm th}=1$. This constant is independent of the degree distribution altogether: the derivation uses only the identity $\sum_k k p_k = \langle k\rangle$, which holds for every degree distribution (Appendix~\ref{app:deriv}, Eq.~\eqref{eq:kappa-infinity}). It may be above or below unity depending on $\bar\phi$, defining a critical threshold
\begin{equation}
    \bar\phi_{\rm crit} = (1-\rho_0)\,\frac{\beta_{\rm th}}{1+\beta_{\rm th}} .
    \label{eq:phi-crit}
\end{equation}
For $\bar\phi<\bar\phi_{\rm crit}$, $\kappa_\infty>1$ and the system can cascade even at arbitrarily high $z$; for $\bar\phi>\bar\phi_{\rm crit}$, $\kappa_\infty<1$ and cascades at large $z$ are then decided by the sign of $1-\alpha_{\rm th}$, growing for $\alpha_{\rm th}<1$ and decaying for $\alpha_{\rm th}>1$.

\begin{figure}[t]
\centering
\fbox{\parbox{0.6\columnwidth}{\centering\vspace{2cm} Placeholder for Fig.~4\vspace{2cm}}}
\caption{Branching ratio $\kappa$ as a function of mean degree $z$ for beta thresholds on ER graphs at various $\alpha_{\rm th}$. Curves are colored by $\alpha_{\rm th}$; the $\kappa=1$ threshold is marked. For $\alpha_{\rm th}<1$, $\kappa$ grows with $z$; for $\alpha_{\rm th}>1$, it decays. At $\alpha_{\rm th}=1$, $\kappa$ approaches the constant $\kappa_\infty$.}
\label{fig:beta-er-kappa}
\end{figure}

\subsection{The re-entrant window and the $\alpha_{\rm th}=1$ transition}

For $\alpha_{\rm th}>1$, the cascade window in $z$ has both an upper and a lower phase transition, as in Watts' original model. As $\bar\phi$ increases, the upper transition shifts leftward and the window shrinks. For sufficiently large $\bar\phi$, the $\kappa(z)$ curve develops a local maximum that dips below unity before decaying to zero, creating a re-entrant window: the system cascades at intermediate $z$, stops at slightly larger $z$, and cascades again at still larger $z$ before stopping for good. This re-entrant window exists only for $\alpha_{\rm th}>1$; at $\alpha_{\rm th}=1$ the $\kappa$ asymptote becomes the constant $\kappa_\infty$, and the window closes.

Numerically, we locate the re-entrant window in $(z,\alpha_{\rm th})$ space by solving $\partial\kappa/\partial z = 0$ at fixed $\bar\phi$. The resulting critical $\alpha_{\rm th}$ at which the window vanishes is always $1$, independent of $\bar\phi$ and $\beta_{\rm th}$. This independence follows directly from the asymptotic form. The exponent $1-\alpha_{\rm th}$ passes through zero at $\alpha_{\rm th}=1$, changing the curvature of $\kappa(z)$.

The $\alpha_{\rm th}=1$ boundary is therefore central to the phase diagram. It separates regimes where cascades eventually decay at large $z$, for $\alpha_{\rm th}>1$, from those where they can persist indefinitely, for $\alpha_{\rm th}<1$. This boundary depends only on the near-zero tail of the threshold distribution, not on its detailed shape, and it holds even when the degree distribution is heterogeneous, as we show below.

\section{Beta Binomial Topology}
\label{sec:bbtopo}

The four limiting cases of \S\ref{sec:fourcases} span the extremes of topological heterogeneity, but they are discrete: the degree distribution is either power law or Poisson, with no way to move continuously between them. To treat degree skewness as a continuous parameter, we replace the degree sequence with a beta-binomial distribution, the discrete analog of the beta family of \S\ref{sec:betath}. This section defines the family and, at constant threshold, isolates how degree skewness alone shapes the cascade window and eventually suppresses cascades at high connectivity. The coupling to a tunable threshold distribution is deferred to \S\ref{sec:phase}. The derivations are given in Appendix~\ref{app:deriv}.

\subsection{Definition and properties}

The beta-binomial distribution is the compound distribution obtained by drawing a success probability $p$ from a beta distribution and then drawing $k$ from $\text{Binomial}(n,p)$. Its probability mass function is
\begin{equation}
    P(k; n, \alpha_d, \beta_d) = \binom{n}{k} \frac{B(k+\alpha_d, n-k+\beta_d)}{B(\alpha_d, \beta_d)},
    \label{eq:bb-pmf}
\end{equation}
where $n$ is the maximum possible degree, $k$ is an integer up to $n$, $\alpha_d$ and $\beta_d$ are shape parameters, and $B(a,b)=\Gamma(a)\Gamma(b)/\Gamma(a+b)$. The mean and variance are
\begin{align}
    \langle k \rangle &= n\,\frac{\alpha_d}{\alpha_d+\beta_d},\\
    \text{Var}(k) &= n\,\frac{\alpha_d\beta_d}{(\alpha_d+\beta_d)^2}\,
    \frac{\alpha_d+\beta_d+n}{\alpha_d+\beta_d+1}.
\end{align}
Relative to a binomial with the same mean, the variance is inflated by a factor $(\alpha_d+\beta_d+n)/(\alpha_d+\beta_d+1)>1$, reflecting the extra heterogeneity from the beta mixing.

\subsection{Continuum from ER to power law}

The shape parameter $\alpha_d$ controls the skewness of the degree distribution. In the limit $\alpha_d\to0^+$ with $\langle k\rangle$ held fixed by taking $\beta_d\to\infty$, the beta-binomial becomes power-law-like: the probability mass at high $k$ decays as $k^{-(1-\alpha_d)}$, up to an exponential cutoff at $k \sim \langle k\rangle/\alpha_d$ (Appendix~\ref{app:deriv}, Eq.~\eqref{eq:pk-asym}). As $\alpha_d$ increases, the tail shortens and the distribution becomes increasingly concentrated around its mean. In the opposite limit, $\alpha_d,\beta_d\to\infty$ with $\beta_d\gg\alpha_d$ and $n\to\infty$ such that $\langle k\rangle$ remains finite, the beta mixing concentrates at its mean and the beta-binomial reduces to a Poisson distribution, the Erd\H{o}s-R\'{e}nyi limit (Appendix~\ref{app:deriv}). The beta-binomial thus spans the full range from strongly skewed power-law to nearly Poisson degree distributions, with $\alpha_d$ a single continuous control parameter.

\subsection{Cascade condition at constant threshold}
\label{sec:bb-const}

Isolating the role of topology requires fixing the thresholds. At constant threshold $\bar\phi$, a node of degree $k$ is vulnerable when $1/k\ge\bar\phi$, i.e.\ $k\le 1/\bar\phi$, and the cascade parameter of Eq.~\eqref{eq:kappa_intro} reduces to
\begin{equation}
    \kappa = (1-\rho_0)\sum_{k=2}^n \frac{k(k-1)P(k; n,\alpha_d,\beta_d)}{\langle k\rangle}\,
    \Theta\!\left(\frac{1}{k}-\bar\phi\right),
    \label{eq:bb-const-kappa}
\end{equation}
where the Heaviside factor truncates the sum at $k_{\max}=\lfloor 1/\bar\phi\rfloor$. The cascade condition is then a function of $\langle k\rangle$ and $\alpha_d$ alone, and the contour $\kappa=1$ in the $(\langle k\rangle,\alpha_d)$ plane separates the cascading and non-cascading regimes, as shown in Fig.~\ref{fig:const-bb-contours}. The constant-threshold case is recovered from the beta thresholds of \S\ref{sec:phase} in the limit $\alpha_{\rm th}\to\infty$.

\begin{figure}[t]
\centering
\fbox{\parbox{0.6\columnwidth}{\centering\vspace{2cm} Placeholder for Fig.~5\vspace{2cm}}}
\caption{Branching ratio $\kappa$ contours at $\kappa=1$ in the $(\langle k\rangle,\alpha_d)$ plane for constant threshold on beta-binomial networks. Contours are shown for several threshold values $\bar\phi$; increasing degree skewness (decreasing $\alpha_d$) moves the onset to lower mean degree and shrinks the cascade window, until for sufficiently small $\alpha_d$ the contour disappears and the network cannot cascade at all.}
\label{fig:const-bb-contours}
\end{figure}

\subsection{Degree skewness suppresses cascades at high connectivity}

At constant threshold the sum in Eq.~\eqref{eq:bb-const-kappa} is finite, truncated at $k_{\max}=\lfloor 1/\bar\phi\rfloor$, so its $z$ dependence enters only through the degree distribution. In the thermodynamic limit $n\to\infty$ at fixed $z=\langle k\rangle$, the small-degree mass of the beta-binomial is $P(k)\approx P_0\,k^{\alpha_d-1}r^k$ with $P_0\propto z^{-\alpha_d}$ and $r\to 1$ (Appendix~\ref{app:deriv}, Eq.~\eqref{eq:pk-asym}), and Eq.~\eqref{eq:bb-const-kappa} gives
\begin{equation}
    \kappa \;\sim\; C_d\, z^{-(\alpha_d+1)} \;\xrightarrow{z\to\infty}\; 0,
    \qquad
    C_d = (1-\rho_0)\,\frac{\Gamma(\alpha_d+\beta_d)}{\Gamma(\beta_d)}
    \,\frac{m_d^{\alpha_d}}{\Gamma(\alpha_d)}
    \sum_{k=2}^{k_{\max}} k(k-1)\,\frac{\Gamma(k+\alpha_d)}{k!},
    \label{eq:const-kappa-asym}
\end{equation}
with $m_d=\alpha_d/(\alpha_d+\beta_d)$; the terms $r^k\approx 1$ on the finite range $k\le k_{\max}$. The decay exponent is set entirely by $\alpha_d$: more skewed degree distributions (smaller $\alpha_d$) decay slower, but every beta-binomial topology eventually suppresses cascades at sufficiently high connectivity. This is the beta-binomial counterpart of the structural bound $\gamma_d\approx3.48$ of \S\ref{sec:pl-gauss}, with the difference that the bound here is asymptotic in $z$ rather than absolute; at any finite $z$ the onset and width of the cascade window follow from the contours of Fig.~\ref{fig:const-bb-contours}.

\section{A Doubly Tunable Model}
\label{sec:phase}

The preceding sections developed two independent parameterizations, the beta distribution for thresholds of \S\ref{sec:betath} and the beta-binomial distribution for degrees of \S\ref{sec:bbtopo}. Each provides a continuous family from strongly heterogeneous to nearly homogeneous, with a single shape parameter controlling the skewness, $\alpha_{\rm th}$ for thresholds and $\alpha_d$ for degrees. The beta-threshold analysis of \S\ref{sec:betath} showed that $\alpha_{\rm th}<1$ eliminates the upper phase transition on ER graphs, and the constant-threshold analysis of \S\ref{sec:bbtopo} showed that degree skewness alone suppresses cascades as $z^{-(\alpha_d+1)}$ at high connectivity. We now combine the two parameterizations into a single model, beta thresholds on beta-binomial networks, and show that their interaction organizes cascades into two asymptotic regimes and, in the $(\alpha_d,\alpha_{\rm th})$ plane, three dynamical regions.

\subsection{The control plane}

The two axes of the phase diagram are the shape parameters $\alpha_{\rm th}$ and $\alpha_d$. On the $\alpha_{\rm th}$ axis, small values put most probability near zero, creating many vulnerable nodes, while large values suppress low thresholds; $\alpha_{\rm th}=1$ separates divergent from finite near-zero density. On the $\alpha_d$ axis, small values give a broad right-skewed degree distribution with many low-degree nodes, while large values give a concentrated, nearly symmetric distribution.

The four limiting cases of \S\ref{sec:fourcases} map to the corners of this plane. At $(\alpha_d\to0^+,\alpha_{\rm th}\to0^+)$ the system is PL--PL, at $(\alpha_d\to\infty,\alpha_{\rm th}\to0^+)$ it is ER--PL, at $(\alpha_d\to0^+,\alpha_{\rm th}\to\infty)$ it is PL--Gaussian, and at $(\alpha_d\to\infty,\alpha_{\rm th}\to\infty)$ it is ER--Gaussian.

\subsection{The two large-degree regimes}

Coupling the beta thresholds of \S\ref{sec:betath} to the beta-binomial topology of \S\ref{sec:bbtopo}, the cascade parameter of Eq.~\eqref{eq:kappa_intro} becomes
\begin{equation}
    \kappa = (1-\rho_0)\sum_{k=2}^n \frac{k(k-1)P(k; n,\alpha_d,\beta_d)}{\langle k\rangle}\,
    I_{m/(k\bar\phi)}(\alpha_{\rm th},\beta_{\rm th}),
    \label{eq:bb-full-kappa}
\end{equation}
with $m=\alpha_{\rm th}/(\alpha_{\rm th}+\beta_{\rm th})$ and $I_u$ the regularized incomplete beta function, which reduces to the constant-threshold case of Eq.~\eqref{eq:bb-const-kappa} as $\alpha_{\rm th}\to\infty$. In the thermodynamic limit $n\to\infty$ at fixed mean degree $z=\langle k\rangle$, the sum in Eq.~\eqref{eq:bb-full-kappa} displays two distinct large-$z$ behaviors, distinguished by which distribution's tail controls the sum. The derivation is in Appendix~\ref{app:deriv}; the results are as follows.

Threshold-dominated, $\alpha_{\rm th}<\alpha_d+2$. The near-zero tail of the threshold distribution dominates the sum, and
\begin{equation}
    \kappa \sim C_1(\alpha_d,\alpha_{\rm th},\beta_{\rm th},\bar\phi)\,
    z^{\,1-\alpha_{\rm th}},
    \label{eq:C1}
\end{equation}
recovering the same $z^{1-\alpha_{\rm th}}$ scaling as the ER case of \S\ref{sec:betath}; the degree skewness enters only through the prefactor $C_1$. For $\alpha_{\rm th}>1$ this branch decays at large $z$, while for $\alpha_{\rm th}<1$ it grows without bound, so at sufficiently high connectivity cascades are unavoidable regardless of the degree distribution. This is the mechanism by which the $\alpha_{\rm th}=1$ boundary of \S\ref{sec:betath} survives degree heterogeneity.

Degree-dominated, $\alpha_{\rm th}>\alpha_d+2$. The threshold tail decays so fast that the degree distribution cuts off the sum, and\begin{equation}
    \kappa \sim C_2(\alpha_d,\alpha_{\rm th},\beta_{\rm th},\bar\phi)\,
    z^{-(\alpha_d+1)},
    \label{eq:C2}
\end{equation}
with the decay exponent set entirely by $\alpha_d$, independent of $\alpha_{\rm th}$; the constant-threshold result of \S\ref{sec:bbtopo}, Eq.~\eqref{eq:const-kappa-asym}, is the special case $\alpha_{\rm th}\to\infty$. In this regime the degree skewness itself determines the large-$z$ decay, and cascades are always suppressed at high connectivity.

The boundary $\alpha_{\rm th}=\alpha_d+2$ between the two regimes is a crossover, not a phase transition: the two asymptotic forms match continuously, up to a logarithmic factor at the boundary itself. The crossover is correspondingly slow, so in numerics the boundary is most reliably located from the scaling exponents rather than from the prefactors (Appendix~\ref{app:deriv}).

The special case $\alpha_{\rm th}=1$ deserves emphasis. Here the threshold tail is exactly proportional to $1/k$, and the cascade parameter approaches the constant $\kappa_\infty$ of Eq.~\eqref{eq:kappa-infinity} \emph{for any degree distribution}: the only identity required is $\sum_k k P(k)=\langle k\rangle$, which holds for every $P(k)$ (Appendix~\ref{app:deriv}, Eq.~\eqref{eq:kappainf-univ}). The horizontal boundary $\alpha_{\rm th}=1$ is therefore the same for every degree distribution, and it is this line, not any property of $\beta_d$, that organizes the phase diagram. There is no analogue of the ``large-$\beta_d$ cancellation'' sometimes claimed for the beta-binomial: the onset of cascades is set by the prefactors $C_1$ and $C_2$, which do depend on $\alpha_d$, and the onset contours of Figs.~\ref{fig:const-bb-contours} and~\ref{fig:bb-phase} are not vertical. The approximate separability of the phase diagram rests on the two-regime structure above, in which each regime's exponent is set by a single parameter, $\alpha_{\rm th}$ or $\alpha_d$ (Appendix~\ref{app:deriv}).

\subsection{Three cascade regimes}

The large-$z$ asymptotics of \S\ref{sec:betath} and the preceding subsection determine the phase diagram in the $(\alpha_d,\alpha_{\rm th})$ plane, which divides into three regions, as shown in Fig.~\ref{fig:bb-phase}.

\subsubsection*{Region I, $\alpha_{\rm th}<1$: cascades are unavoidable}

For $\alpha_{\rm th}<1$, the near-zero tail of the threshold distribution is divergent, with $f(x)\sim x^{\alpha_{\rm th}-1}$ and $\alpha_{\rm th}-1<0$. This produces a nonvanishing fraction of nodes with arbitrarily low thresholds, so some nodes are vulnerable regardless of degree. Since $\alpha_{\rm th}<1<\alpha_d+2$ always, the system lies in the threshold-dominated regime, and
\begin{equation}
    \kappa \;\sim\; C_1\;z^{\,1-\alpha_{\rm th}} \;\xrightarrow{z\to\infty}\; \infty ,
\end{equation}
for any $\alpha_d>0$. Cascades are therefore unavoidable at sufficiently high connectivity: for any degree distribution and any mean threshold $\bar\phi$, there is a critical mean degree $z_c$ above which $\kappa>1$. The degree distribution enters only through the prefactor $C_1$, which sets $z_c$ but cannot prevent cascades. This result is independent of $\beta_{\rm th}$ and $\beta_d$: the divergent threshold tail at $x=0$ dominates the sum over degrees, and no amount of topological homogeneity compensates. Region I occupies the entire lower half of the phase diagram, $\alpha_{\rm th}<1$.

\subsubsection*{Region II, $1<\alpha_{\rm th}<\alpha_d+2$: cascades are conditional}

For $\alpha_{\rm th}>1$, the threshold density is finite at $x=0$, and the fraction of vulnerable nodes decays with degree as $F(1/k)\sim k^{-\alpha_{\rm th}}$. The system remains in the threshold-dominated regime, and the cascade parameter decays at large $z$,
\begin{equation}
    \kappa \;\sim\; C_1\;z^{\,1-\alpha_{\rm th}} \;\xrightarrow{z\to\infty}\; 0 ,
\end{equation}
so cascades are suppressed at very high connectivity. However, for intermediate $z$, the $\kappa(z)$ curve may still exceed unity, producing a finite cascade window of the Watts type. The existence and extent of this window depend on the prefactor $C_1$ in Eq.~\eqref{eq:C1}, which encodes the threshold shape $(\alpha_{\rm th},\beta_{\rm th})$, degree shape $(\alpha_d,\beta_d)$, and mean threshold $\bar\phi$. The two boundaries of this region are where the threshold tail stops controlling the dynamics: at $\alpha_{\rm th}=1$ the threshold tail is exactly $1/k$ and the large-$z$ limit becomes the constant $\kappa_\infty$ of Eq.~\eqref{eq:kappa-infinity}; at $\alpha_{\rm th}=\alpha_d+2$ the degree tail takes over entirely (Region III).

\subsubsection*{Region III, $\alpha_{\rm th}>\alpha_d+2$: cascades are degree-dominated}

When the threshold tail decays faster than the degree distribution can supply the high-degree nodes that drive propagation, the sum over $k$ is controlled by the degree distribution, and
\begin{equation}
    \kappa \;\sim\; C_2\;z^{-(\alpha_d+1)} \;\xrightarrow{z\to\infty}\; 0 ,
\end{equation}
with the decay exponent set entirely by $\alpha_d$, independent of $\alpha_{\rm th}$. Cascades are suppressed at high connectivity regardless of how skewed the threshold distribution is. The constant-threshold limit, $\alpha_{\rm th}\to\infty$, always lies in this region: with a constant threshold only nodes of degree below $1/\bar\phi$ are ever vulnerable, so the sum is finite and its large-$z$ decay is set by the degree distribution alone. This recovers the constant-threshold mechanism of \S\ref{sec:bbtopo}, Eq.~\eqref{eq:const-kappa-asym}, and parallels the structural bound $\gamma_d\approx3.48$ of \S\ref{sec:pl-gauss}.

\subsection{The $\alpha_{\rm th}=1$ critical boundary}

The horizontal line $\alpha_{\rm th}=1$ is the main boundary of the phase diagram. It separates the regime where cascades persist at arbitrarily high connectivity, Region I, from the regime where they are eventually suppressed, Regions II and III. Three properties fix this boundary. First, it is shape-independent: the boundary depends only on the near-zero behavior of the threshold distribution, $f\sim x^{\alpha_{\rm th}-1}$, not on its detailed shape, so changing $\beta_{\rm th}$ or $\bar\phi$ shifts the onset of cascades but does not move the boundary. Second, it is topology-independent: at $\alpha_{\rm th}=1$ the cascade parameter approaches the constant $\kappa_\infty$ of Eq.~\eqref{eq:kappa-infinity}, which follows from the single identity $\sum_k k p_k=\langle k\rangle$ and is therefore identical for every degree distribution (Appendix~\ref{app:deriv}, Eq.~\eqref{eq:kappainf-univ}). Third, the re-entrant window vanishes at this boundary. For $\alpha_{\rm th}>1$, the $\kappa(z)$ curve develops a local maximum that can produce a re-entrant cascade window, but at $\alpha_{\rm th}=1$ the large-$z$ asymptote becomes the constant $\kappa_\infty$ and the window closes continuously.

\subsection{Mapping the four limiting cases}

The four corners of the $(\alpha_d,\alpha_{\rm th})$ plane recover the limiting cases of \S\ref{sec:fourcases}.

At $(\alpha_d\to0,\,\alpha_{\rm th}\to0)$, the PL--PL corner, both distributions are maximally skewed. The system falls in Region~I, where cascades are unavoidable at any $\bar\phi$ and any connectivity, consistent with \S\ref{sec:pl-pl}.

At $(\alpha_d\to\infty,\,\alpha_{\rm th}\to0)$, the ER--PL corner, the degree distribution is Poisson and concentrated while thresholds are still heavily skewed. The system remains in Region~I, as found in \S\ref{sec:er-pl}.

At $(\alpha_d\to0,\,\alpha_{\rm th}\to\infty)$, the PL--Gaussian corner, the degree distribution is broad but thresholds are narrow. The system lies in Region~III when $\alpha_d$ is small enough that $\alpha_{\rm th}\ge\alpha_d+2$, and the cascade is degree-limited. For PL degree distributions with $\gamma_d\gtrsim3.48$, the structural bound prevents cascades regardless of $\bar\phi$, as shown in \S\ref{sec:pl-gauss}. In the beta-binomial framework, small $\alpha_d$ produces a similarly broad degree distribution, and the $z^{-(\alpha_d+1)}$ decay ensures suppression.

At $(\alpha_d\to\infty,\,\alpha_{\rm th}\to\infty)$, the ER--Gaussian corner, both distributions are concentrated. The system is in Region~II, and the cascade window is the familiar intermediate range of $z$ from Watts' original model, as described in \S\ref{sec:er-gauss}. The branching ratio $\langle k(k-1)\rangle/\langle k\rangle = z$ grows without bound, so the window is controlled entirely by the threshold.

\subsection{Discussion}

The $(\alpha_d,\alpha_{\rm th})$ phase diagram unifies the four limiting cases into a single continuous framework. Threshold skewness dominates. The $\alpha_{\rm th}=1$ boundary divides the phase diagram into a lower half where cascades are always possible, Region~I, and an upper half where they are conditional, Regions II and III. The boundary depends only on the threshold distribution, not on the degree distribution, because at $\alpha_{\rm th}=1$ the cascade parameter is set by the degree mean alone. Within the conditional half, the diagonal line $\alpha_{\rm th}=\alpha_d+2$ is a crossover, not a phase transition: the two asymptotic regimes match continuously and the cascade window shrinks smoothly to zero. Its role is to mark where the binding constraint changes. Above it, the threshold tail is thin and the degree distribution controls the decay exponent; below it, the threshold tail controls both the exponent and the shape of the window, with the degree distribution modulating only the onset $z_c$ through the prefactor $C_1$: broader degree distributions (smaller $\alpha_d$) lower $z_c$.

The beta-binomial interpolates between skewed and symmetric degree distributions. Unlike power-law degree distributions it has finite moments of all orders, yet the $z^{-(\alpha_d+1)}$ decay in Region~III produces an effective bound for small $\alpha_d$ that parallels the $\gamma_d\approx3.48$ bound from power-law networks: both are the statement that a sufficiently fat degree tail starves the cascade of the intermediate-degree nodes that sustain propagation.

These results establish that the interplay of threshold and degree skewness is the fundamental mechanism governing cascade dynamics. The beta--beta-binomial parameterization is a tractable analytical framework for this interplay, and the $(\alpha_d,\alpha_{\rm th})$ phase diagram summarizes when cascades are possible, when they are unavoidable, and when topology or the threshold distribution is the binding constraint.

\begin{figure}[t]
\centering
\fbox{\parbox{0.6\columnwidth}{\centering\vspace{2cm} Placeholder for Fig.~6\vspace{2cm}}}
\caption{Branching ratio $\kappa$ contours at $\kappa=1$ on the left and final cascade fraction $\rho$ on the right for the beta--beta-binomial model at fixed $\bar\phi$. The $\kappa$ contours in the $(\alpha_d,\langle k\rangle)$ plane show the three dynamical regimes, separated by the horizontal boundary $\alpha_{\rm th}=1$ and the diagonal crossover $\alpha_{\rm th}=\alpha_d+2$. The $\rho$ panels show the corresponding bifurcation diagrams for representative $\alpha_{\rm th},\beta_{\rm th}$ pairs.}
\label{fig:bb-phase}
\end{figure}

\appendix
\section{The beta and beta-binomial: derivations and asymptotics}
\label{app:deriv}

This appendix collects the derivations behind the asymptotic results quoted in Sections \ref{sec:betath}--\ref{sec:phase}. Throughout, we use the self-flipping convention of \citet{GleesonCahalane2007} and the cascade parameter of Eq.~\eqref{eq:kappa_intro}. The seed prefactor $(1-\rho_0)$ is carried explicitly.

\subsection{The rescaled beta threshold and its near-zero tail}

The beta distribution has density
\begin{equation}
    f(x;\alpha,\beta) = \frac{x^{\alpha-1}(1-x)^{\beta-1}}{B(\alpha,\beta)},
\end{equation}
and near $x=0$ behaves as $f(x) = x^{\alpha-1}/B(\alpha,\beta)\,[1+O(x)]$; the exponent $\alpha-1$ encodes the near-zero density. The rescaled threshold $\Phi=(\bar\phi/m)X$ with $m=\alpha/(\alpha+\beta)$ has CDF
\begin{equation}
    F(x) = \Pr\left(X \le \frac{m x}{\bar\phi}\right)
    = \min\!\left(1,\; I_{m x/\bar\phi}(\alpha,\beta)\right),
\end{equation}
where $I_u$ is the regularized incomplete beta function; the truncation accounts for rescaled thresholds above unity, which are never vulnerable.

The regularized incomplete beta function has the small-argument expansion obtained by expanding $(1-t)^{\beta-1}$ under the integral,
\begin{equation}
    I_u(\alpha,\beta)
    = \frac{u^\alpha}{\alpha\,B(\alpha,\beta)}
    \left[1 + \frac{1-\beta}{1+\alpha}\,u + O(u^2)\right],
    \label{eq:Iu-asym}
\end{equation}
valid as $u\to0$. Evaluated at the cascade condition $x=1/k$, the argument is $u = m/(k\bar\phi)$, and
\begin{equation}
    F\!\left(\frac{1}{k}\right)
    = C_{\rm th}\, k^{-\alpha_{\rm th}}\,[1+O(1/k)],
    \qquad
    C_{\rm th} = \frac{m^{\alpha_{\rm th}}}{\alpha_{\rm th}\,B(\alpha_{\rm th},\beta_{\rm th})\,\bar\phi^{\alpha_{\rm th}}}.
    \label{eq:Ftail}
\end{equation}
Only the exponent $\alpha_{\rm th}$ survives in the large-$k$ limit; all other details of the threshold shape enter through the constant $C_{\rm th}$. The boundary at $\alpha_{\rm th}=1$ is therefore set by the near-zero exponent alone.

\subsection{The cascade parameter on Erd\H{o}s--R\'{e}nyi graphs}

Substituting the Poisson degree distribution $p_k = e^{-z}z^k/k!$ with $\langle k\rangle=z$ into Eq.~\eqref{eq:kappa_intro} gives
\begin{equation}
    \kappa(z)
    = (1-\rho_0)\sum_{k=2}^{\infty}
    \frac{k(k-1)}{z}\,\frac{e^{-z}z^k}{k!}\,I_{m/(k\bar\phi)}(\alpha_{\rm th},\beta_{\rm th})
    = (1-\rho_0)\sum_{k=2}^{\infty}
    \frac{e^{-z}z^{k-1}}{(k-2)!}\,I_{m/(k\bar\phi)}(\alpha_{\rm th},\beta_{\rm th}),
\end{equation}
where we used $k(k-1)/k! = 1/(k-2)!$.

For $z\to\infty$ the weight $\sum_{k\ge2} e^{-z}z^{k-1}/(k-2)! = z$ is concentrated at $k\approx z$. Substituting Eq.~\eqref{eq:Ftail},
\begin{equation}
    \kappa(z) \approx (1-\rho_0)\,C_{\rm th}\,
    \sum_{k=2}^{\infty} \frac{e^{-z}z^{k-1}}{(k-2)!}\,k^{-\alpha_{\rm th}}.
\end{equation}
Writing $k^{-\alpha_{\rm th}} = z^{-\alpha_{\rm th}}(k/z)^{-\alpha_{\rm th}}$ and expanding the last factor about $k=z$ gives $(k/z)^{-\alpha_{\rm th}} = 1 - \alpha_{\rm th}(k-z)/z + O(1/z^2)$; the linear term vanishes by the mean identity $\sum_k (k-z)\,p_k^{\rm SB}=0$ for the size-biased Poisson weight, leaving
\begin{equation}
    \kappa \sim (1-\rho_0)\,C_{\rm th}\,z^{1-\alpha_{\rm th}}
    = (1-\rho_0)\,\frac{m^{\alpha_{\rm th}}}{\alpha_{\rm th} B(\alpha_{\rm th},\beta_{\rm th})}
    \,\bar\phi^{-\alpha_{\rm th}}\,z^{1-\alpha_{\rm th}} .
\end{equation}
The sign of $1-\alpha_{\rm th}$ fixes the large-$z$ fate of the cascade: growth for $\alpha_{\rm th}<1$, decay for $\alpha_{\rm th}>1$.

At the boundary $\alpha_{\rm th}=1$ the CDF is elementary, $I_u(1,\beta) = 1-(1-u)^\beta = \beta u + O(u^2)$, so $F(1/k) = \beta_{\rm th} m/(\bar\phi\,k) + O(k^{-2})$. The cascade parameter becomes
\begin{equation}
    \kappa(z) \approx (1-\rho_0)\,\frac{\beta_{\rm th} m}{\bar\phi}
    \sum_{k=2}^{\infty}\frac{e^{-z}z^{k-1}}{(k-2)!}\,\frac{1}{k}
    = (1-\rho_0)\,\frac{\beta_{\rm th} m}{\bar\phi}\left(1 - \frac{1}{z} + O(e^{-z})\right),
\end{equation}
and so
\begin{equation}
    \kappa_\infty = \lim_{z\to\infty}\kappa(z)
    = (1-\rho_0)\,\frac{\beta_{\rm th}\,m}{\bar\phi}
    = (1-\rho_0)\,\frac{\beta_{\rm th}}{(1+\beta_{\rm th})\,\bar\phi},
    \label{eq:kappainf-univ}
\end{equation}
where the last equality uses $B(1,\beta_{\rm th})=1/\beta_{\rm th}$ and $m=1/(1+\beta_{\rm th})$. Setting $\kappa_\infty=1$ gives the critical mean threshold
\begin{equation}
    \bar\phi_{\rm crit} = (1-\rho_0)\,\frac{\beta_{\rm th}}{1+\beta_{\rm th}} .
\end{equation}

\subsection{The beta-binomial degree distribution in the thermodynamic limit}

The beta-binomial PMF of Eq.~\eqref{eq:bb-pmf} satisfies the exact ratio recurrence
\begin{equation}
    \frac{P(k+1;n,\alpha_d,\beta_d)}{P(k;n,\alpha_d,\beta_d)}
    = \frac{(n-k)(k+\alpha_d)}{(k+1)(n-k-1+\beta_d)},
    \label{eq:bb-ratio}
\end{equation}
obtained from the beta-function ratio $B(k+1+\alpha_d,n-k-1+\beta_d)/B(k+\alpha_d,n-k+\beta_d) = (k+\alpha_d)/(n-k-1+\beta_d)$ and the binomial ratio $\binom{n}{k+1}/\binom{n}{k} = (n-k)/(k+1)$.

We work in the thermodynamic limit $n\to\infty$ with $\alpha_d,\beta_d$ fixed and mean degree $z = n m_d$, $m_d=\alpha_d/(\alpha_d+\beta_d)$, held fixed. For $k\ll n$, the first factor of Eq.~\eqref{eq:bb-ratio} is $(n-k)/(n-k-1+\beta_d) \approx n/(n+\beta_d) \equiv r$, so
\begin{equation}
    P(k;n,\alpha_d,\beta_d) \approx P(0)\, \frac{\Gamma(k+\alpha_d)}{\Gamma(\alpha_d)\,k!}\,r^k,
    \label{eq:pk-rec}
\end{equation}
with $r = n/(n+\beta_d)\to 1$. The normalization constant is
\begin{equation}
    P(0) = \frac{B(\alpha_d,\,n+\beta_d)}{B(\alpha_d,\beta_d)}
    = \frac{\Gamma(\alpha_d+\beta_d)}{\Gamma(\beta_d)}\,
    \frac{\Gamma(n+\beta_d)}{\Gamma(n+\alpha_d+\beta_d)}
    \approx \frac{\Gamma(\alpha_d+\beta_d)}{\Gamma(\beta_d)}\,
    m_d^{\alpha_d}\,z^{-\alpha_d},
\end{equation}
using $\Gamma(n+\beta_d)/\Gamma(n+\alpha_d+\beta_d) \approx (n+\beta_d)^{-\alpha_d}$ and $n^{-1}=m_d/z$ for large $n$.

For $k\gg1$ within $k\ll n$, Stirling's formula $\Gamma(k+\alpha_d)/k! \sim k^{\alpha_d-1}$ turns Eq.~\eqref{eq:pk-rec} into
\begin{equation}
    P(k) \approx P_0\, k^{\alpha_d-1}\,r^k,
    \qquad
    P_0 = \frac{\Gamma(\alpha_d+\beta_d)}{\Gamma(\beta_d)}\,
    \frac{m_d^{\alpha_d}}{\Gamma(\alpha_d)}\,z^{-\alpha_d},
    \label{eq:pk-asym}
\end{equation}
a power-law tail $P(k)\propto k^{-(1-\alpha_d)}$ with an exponential cutoff set by $r^k = e^{-k/k_c}$ at
\begin{equation}
    k_c = \frac{1}{1-r} = \frac{n+\beta_d}{\beta_d}
    = \frac{z(\alpha_d+\beta_d)}{\alpha_d\beta_d} + 1
    \;\sim\; \frac{z}{\alpha_d}
    \qquad(\beta_d\gg\alpha_d).
\end{equation}
The cutoff is the footprint of the finite maximum degree $n$; at fixed mean $z$ it scales as $z/\alpha_d$, moving to infinity only as $\alpha_d\to0$. This is the precise sense in which $\alpha_d\to0^+$ with $\beta_d\to\infty$ gives a power-law-like degree distribution, and in which increasing $\alpha_d$ concentrates the distribution.

The ER limit requires the beta mixing to concentrate, i.e.\ $\alpha_d,\beta_d\to\infty$ with $p_0=\alpha_d/(\alpha_d+\beta_d)$ fixed; the beta-binomial then tends to $\mathrm{Bin}(n,p_0)$, and for $n\to\infty$ with $z=np_0$ fixed, $\mathrm{Bin}(n,p_0)\to\mathrm{Poisson}(z)$ by the Poisson limit theorem. Note that $\alpha_d=\beta_d\to\infty$ does \emph{not} give this limit: it gives $p_0=1/2$, a binomial concentrated at $n/2$, whose mean diverges. The homogeneous (ER) corner of the phase diagram therefore corresponds to $\beta_d\gg\alpha_d$ with both large, as used in the numerics.

\subsection{The cascade parameter on beta-binomial graphs}

Substituting Eq.~\eqref{eq:pk-rec} and Eq.~\eqref{eq:Ftail} into Eq.~\eqref{eq:kappa_intro} with $\langle k\rangle=z$ gives
\begin{equation}
    \kappa \approx (1-\rho_0)\,\frac{C_{\rm th}\,P(0)}{z\,\Gamma(\alpha_d)}
    \sum_{k=2}^{n}\frac{\Gamma(k+\alpha_d)}{k!}\,k(k-1)\,k^{-\alpha_{\rm th}}\,r^k .
\end{equation}
With $k(k-1)/k! = 1/(k-2)!$ and $j=k-2$,
\begin{equation}
    \kappa \approx (1-\rho_0)\,\frac{C_{\rm th}\,P(0)}{z\,\Gamma(\alpha_d)}
    \sum_{j=0}^{\infty} \frac{\Gamma(j+2+\alpha_d)}{j!}\,(j+2)^{-\alpha_{\rm th}}\,r^{j+2},
    \label{eq:bb-kappa-asym}
\end{equation}
where the upper limit is extended to infinity because $r<1$ makes the tail negligible. The asymptotic behavior as $z\to\infty$ (so $r\to1$) is governed by the tail of the summand, which behaves as
\begin{equation}
    \frac{\Gamma(j+2+\alpha_d)}{j!}\,(j+2)^{-\alpha_{\rm th}}
    \;\sim\; j^{\,1+\alpha_d-\alpha_{\rm th}}
    \qquad (j\to\infty),
\end{equation}
by Stirling's formula. Two cases arise, separated by whether this exponent exceeds or falls below $-1$.

Threshold-dominated, $\alpha_{\rm th}<\alpha_d+2$. The summand decays slower than $1/j$; the sum diverges as $r\to1$ and is cut off at $j\sim k_c = 1/(1-r)$. Replacing the sum by an integral,
\begin{equation}
    \sum_{j=0}^{\infty} \frac{\Gamma(j+2+\alpha_d)}{j!}\,(j+2)^{-\alpha_{\rm th}}\,r^{j+2}
    \;\sim\; \int_0^{k_c} dj\; j^{\,1+\alpha_d-\alpha_{\rm th}}
    = \frac{k_c^{2+\alpha_d-\alpha_{\rm th}}}{2+\alpha_d-\alpha_{\rm th}}
    \;\propto\; z^{\,2+\alpha_d-\alpha_{\rm th}},
\end{equation}
and assembling $P(0)$ and $C_{\rm th}$,
\begin{equation}
    \kappa \sim C_1\, z^{\,1-\alpha_{\rm th}},
    \qquad
    C_1 = (1-\rho_0)\,C_{\rm th}\,
    \frac{\Gamma(\alpha_d+\beta_d)}{\Gamma(\beta_d)}\,
    \frac{m_d^{\alpha_d}}{\Gamma(\alpha_d)}\,
    \left(\frac{\alpha_d+\beta_d}{\alpha_d\beta_d}\right)^{2+\alpha_d-\alpha_{\rm th}}
    \frac{1}{2+\alpha_d-\alpha_{\rm th}} .
\end{equation}
The large-$z$ exponent, $1-\alpha_{\rm th}$, is identical to the ER result, Eq.~\eqref{eq:er-asym}; the degree distribution affects only the prefactor. This is the threshold-dominated regime.

Degree-dominated, $\alpha_{\rm th}>\alpha_d+2$. The summand decays faster than $1/j$; the sum converges as $r\to1$ to a constant,
\begin{equation}
    S(\alpha_d,\alpha_{\rm th}) \equiv
    \sum_{j=0}^{\infty} \frac{\Gamma(j+2+\alpha_d)}{j!}\,(j+2)^{-\alpha_{\rm th}} < \infty ,
\end{equation}
and
\begin{equation}
    \kappa \sim C_2\, z^{-(\alpha_d+1)},
    \qquad
    C_2 = (1-\rho_0)\,C_{\rm th}\,
    \frac{\Gamma(\alpha_d+\beta_d)}{\Gamma(\beta_d)}\,
    \frac{m_d^{\alpha_d}}{\Gamma(\alpha_d)}\,
    S(\alpha_d,\alpha_{\rm th}) .
\end{equation}
The decay exponent, $-(\alpha_d+1)$, is set entirely by the degree distribution; the threshold enters only through the prefactor. This is the degree-dominated regime.

Boundary, $\alpha_{\rm th}=\alpha_d+2$. The summand behaves as $1/j$, the sum diverges logarithmically, $S\sim\ln k_c\sim\ln z$, and
\begin{equation}
    \kappa \sim C\,z^{-(\alpha_d+1)}\ln z .
\end{equation}
The crossover between the two regimes is therefore slow, which is why the boundary is most reliably identified numerically from the scaling exponents rather than from the prefactors. This logarithmic crossover is confirmed by direct evaluation of the sum.

Two remarks on the prefactors. First, Eq.~\eqref{eq:Ftail} is only the leading small-argument behavior of the beta CDF; for small $\alpha_{\rm th}$ the low-degree terms, where the argument $m/(k\bar\phi)$ is not small, contribute at an order-one level, so Eqs.~\eqref{eq:C1} and \eqref{eq:C2} give the correct scaling but the constants are best obtained numerically. The exponents, which fix the phase boundaries, are exact. In the degree-dominated regime the prefactor also receives an order-one contribution from the low degrees with $m/(k\bar\phi)\ge1$, where the CDF saturates at unity rather than following Eq.~\eqref{eq:Ftail}; as $\alpha_{\rm th}\to\infty$ these terms dominate, and the limit recovers the constant-threshold prefactor $C_d$ of Eq.~\eqref{eq:const-kappa-asym}. Second, there is no ``large-$\beta_d$ cancellation'' of the kind that would make the $\kappa=1$ contours vertical: the ratio appearing in Eq.~\eqref{eq:C1}, $R = [(\alpha_d+\beta_d)/\alpha_d]^{2-\alpha_{\rm th}}\,B(\alpha_d+2-\alpha_{\rm th},\beta_d)/B(\alpha_d,\beta_d)$, converges as $\beta_d\to\infty$, but to $\Gamma(\alpha_d+2-\alpha_{\rm th})/[\Gamma(\alpha_d)\,\alpha_d^{2-\alpha_{\rm th}}]$, not to unity, and the onset $z_c = C_1^{1/(\alpha_{\rm th}-1)}$ in the threshold-dominated regime depends on $\alpha_d$ through $C_1$. The approximate separability of the phase diagram rests on the two-regime structure itself: in each regime the large-$z$ exponent is set by a single parameter, $\alpha_{\rm th}$ in the threshold-dominated regime and $\alpha_d$ in the degree-dominated regime, and the boundary between the regimes is the line $\alpha_{\rm th}=\alpha_d+2$.

\subsection{The $\alpha_{\rm th}=1$ boundary}

At $\alpha_{\rm th}=1$, Eq.~\eqref{eq:Ftail} is exact to leading order in $1/k$, $F(1/k) = \beta_{\rm th} m/(\bar\phi\,k)\,[1+O(1/k)]$, for \emph{any} degree distribution. Substituting into Eq.~\eqref{eq:kappa_intro},
\begin{equation}
    \kappa
    \approx (1-\rho_0)\,\frac{\beta_{\rm th} m}{\bar\phi}
    \sum_{k}\frac{k(k-1)}{\langle k\rangle}\,p_k\,\frac{1}{k}
    = (1-\rho_0)\,\frac{\beta_{\rm th} m}{\bar\phi}
    \,\frac{\langle k\rangle - 1}{\langle k\rangle},
\end{equation}
which tends to $\kappa_\infty = (1-\rho_0)\beta_{\rm th}m/\bar\phi$ as $\langle k\rangle\to\infty$. The only input is the definition of the mean degree, $\sum_k k p_k = \langle k\rangle$, so $\kappa_\infty$ is identical for every degree distribution. The boundary is therefore the same for every degree distribution: at $\alpha_{\rm th}=1$ the large-$z$ fate of the cascade is set by the mean connectivity alone, and the horizontal line $\alpha_{\rm th}=1$ is fixed in the phase diagram for all $\alpha_d$.

\bibliography{references}

\end{document}
